\section{A manifold criterion for specialness and fundamental groups}

Here $X$ may have any positive dimension $n$, and $D$ is still a nonempty smooth connected reduced divisor. Write $\cX=\sqrt[2]{(X,D)}$.

If a map $g:Y\to X$ from a smooth manifold satisfies $g^*D=2E$, then the divisor line $\OO_Y(E)$ and its canonical section, with their squared identification with $g^*(\OO_X(D),s_D)$, define a map
\[
 f:Y\longrightarrow\cX,\qquad \pi f=g.
\]
Equivalently, a local defining equation of $D$ pulls back to a square after choosing a local square root of a unit. Choices differ by sign and give precisely the root-chart transition rule. Orbifold bundles and differential maps therefore pull back to ordinary ones on $Y$.

\begin{proposition}[Transfer criterion]\label{prop:transfer}
Let $g:Y\to X$ be a surjective morphism with $Y$ connected smooth projective. Assume:
\begin{enumerate}
\item\label{item:even} $g^*D=2E$ for an effective integral divisor $E$ on $Y$;
\item\label{item:fibres} over a nonempty Zariski open $B\subset X\setminus D$, the map $g$ is smooth with connected fibres $F$, and $\Omega_F^1$ is trivial;
\item\label{item:ordinary} above a general point of each prime divisor $P\ne D$ of $X$, there is a point at which $g$ is a submersion;
\item\label{item:root} above a general point of $D$, there is a point at which a local lift to the root chart is a submersion.
\end{enumerate}
If $\cX$ is special, then $Y$ is special. There is also a surjection
\begin{equation}\label{eq:group-surjection}
 \pi_1(Y)\twoheadrightarrow
 \pi_1(X\setminus D)/\langle\!\langle\gamma^2\rangle\!\rangle.
\end{equation}
Consequently Assumption~\ref{ass:manifold} implies virtual abelianity of the group on the right.
\end{proposition}

\begin{proof}[Specialness]
The map to the root orbifold described above will be denoted by $f$. Suppose that $i:L\to\Omega_Y^p$ contradicts specialness of $Y$. We first descend the plurisection ratios, then the associated differential tensors, and finally test their poles over base divisors. Choose sections $s_0,\ldots,s_p$ of $L^{\otimes m}$ as in Lemma~\ref{lem:ratios}, and put $u_a=s_a/s_0$.

For a smooth general fibre $F$ over $B$, the exact cotangent sequence
\[
 0\longrightarrow\OO_F^{\oplus n}
 \longrightarrow\Omega_Y^1|_F
 \longrightarrow\Omega_F^1\longrightarrow0
\]
gives a filtration of $\Omega_Y^p|_F$ by subbundles with trivial graded pieces. Choose $F$ for which both $i|_F$ and $s_0|_F$ are nonzero. The first nonzero projection of $i|_F$ to a graded piece, followed by a suitable coordinate projection, gives a nonzero section $t$ of $L^{-1}|_F$. Then
\[
 t^m s_0|_F\in H^0(F,\OO_F)
\]
is nonzero. It is a nonzero constant because $F$ is compact and connected. In particular $t$ and $s_0|_F$ have no zeros, and every $u_a$ is constant on $F$. Lemma~\ref{lem:rational-descent} gives
\[
 u_a=g^*v_a,\qquad v_a\in\CC(X).
\]
Their independent differentials imply $p\le n$ and
\[
 \sigma=dv_1\wedge\cdots\wedge dv_p\ne0.
\]
By Lemma~\ref{lem:ratios}, the generic image of $L$ is the line of $g^*\sigma$.

We next descend the plurisections through their images as differential tensors. Write
\begin{equation}\label{eq:tensor-coefficients}
 i^{\otimes m}(s_a)=h_a(g^*\sigma)^{\otimes m},
 \qquad h_a\in\CC(Y).
\end{equation}
Shrink the smooth base open so that $\sigma$ is holomorphic and nowhere zero. Because $g$ is a submersion there, $g^*\sigma$ is nowhere zero at every point of the full inverse image. Equation~\eqref{eq:tensor-coefficients} makes each $h_a$ holomorphic there: locally one divides by a nonvanishing coefficient of the reference tensor. Compact connected fibres force $h_a$ to be constant on each such fibre. Therefore
\[
 h_a=g^*b_a\quad(b_a\in\CC(X)),\qquad b_a/b_0=v_a.
\]

Let $\tau=\pi^*\sigma$ be the meromorphic form on the smooth root charts. Lemma~\ref{lem:saturation} gives an orbifold line bundle $K\to\Omega_{\cX}^p$ with generic direction $\tau$. We claim that every tensor
\begin{equation}\label{eq:descended-tensors}
 (\pi^*b_a)\tau^{\otimes m}
\end{equation}
is holomorphic. Its differential pullback to $Y$ is holomorphic by \eqref{eq:tensor-coefficients}. At the general point of any ordinary base prime, condition~\ref{item:ordinary} supplies a local section of a submersive lift. At $D$, condition~\ref{item:root} does the same in the root chart. Pulling back along such a local section proves that \eqref{eq:descended-tensors} has no pole at that divisorial point. Any polar chart divisor belongs to the invariant global polar locus of this rational tensor. Its finite image is a base divisor, so equivariance and the preceding tests exclude every polar component. Extension across codimension two proves the claim everywhere.

By Lemma~\ref{lem:saturation}, the tensors in \eqref{eq:descended-tensors} are sections of $K^{\otimes m}$. Their ratios are the independent $v_a$, so $\kappa(K)\ge p$. Conversely, $f^*K\to\Omega_Y^p$ is a nonzero line map: on the dense ordinary open it is the differential pullback by a dominant map. Pullback preserves independence of section ratios, and Lemma~\ref{lem:ratios} gives
\[
 \kappa(K)\le\kappa(f^*K)\le p.
\]
Thus $K$ contradicts specialness of $\cX$. This proves specialness of $Y$.
\end{proof}

\begin{proof}[Fundamental groups and conclusion]
Set $Y^\circ=Y\setminus g^{-1}(D)$ and $X^\circ=X\setminus D$. For a smooth manifold and a divisor, the inclusion of the complement surjects on fundamental groups, with kernel normally generated by meridians of the irreducible divisor components. Indeed, perturb loops off the divisor and perturb null-homotopy discs transversely to its smooth strata, avoiding the strata of complex codimension at least two. Puncturing a disc at its transverse intersections gives the normal generators. This applies even if the reduced support of $g^*D$ is singular.

Consider $\pi_1(Y^\circ)\to\pi_1(X^\circ)\to G$. At the general point of a component of $g^*D$, its meridian maps to a conjugate of $\gamma^{2e}$, where $2e$ is that component's multiplicity. The winding calculation uses a smooth local equation for $D$ and remains valid when the component maps into a smaller subset of $D$. Positive meridians of $D$ are conjugate because $D$ is smooth and connected. Thus all the meridians killed on filling $Y^\circ$ die in $G$, and the map factors through $\pi_1(Y)$.

Over $B$ the map is proper and smooth, hence a differentiable fibre bundle \cite{Ehresmann}. Connectedness of its fibres gives $\pi_1(g^{-1}B)\twoheadrightarrow\pi_1(B)$. General position gives $\pi_1(B)\twoheadrightarrow\pi_1(X^\circ)$. The map just factored is therefore surjective, proving \eqref{eq:group-surjection}. Finally, $Y$ is a connected smooth projective special manifold. Assumption~\ref{ass:manifold} applies, and the image of a finite-index abelian subgroup of $\pi_1(Y)$ is finite-index abelian in $G$.
\end{proof}
